\documentclass[10pt,a4paper]{amsart}
\usepackage[margin=19mm]{geometry}
\usepackage{amsmath,amssymb,mathtools}
\usepackage[scr=rsfs,cal=euler]{mathalfa}
\usepackage{xfrac}
\usepackage[unicode,hidelinks,bookmarks=false]{hyperref}
\newcommand{\ie}{i.e.\ }
\newcommand{\cf}{cf.\ }
\newcommand{\resp}{respectively\ }
\newcommand{\Subsection}{\subsection}
\providecommand{\nobreakdash}{\nobreak\hbox{-}}
\setlength{\emergencystretch}{2em}
\multlinegap0pt
\usepackage{tikz}
\usetikzlibrary{cd}
\usepackage{natbib}
\usepackage{amssymb}
\usepackage{mathtools}
\usepackage{bm}
\usepackage[shortlabels]{enumitem}
\newtheorem{proposition}{Proposition}
\title{Bloch--Suslin Complex and Strong $\mathbb{A}^{1}-$invariance}
\author{Saheb Mohapatra}
\date{Source: arXiv:2512.17118v1 (CC BY 4.0)}
\begin{document}
\maketitle

\noindent\emph{Source and licence.} Bloch--Suslin Complex and Strong $\mathbb{A}^{1}-$invariance. Version arXiv:2512.17118v1 (CC BY 4.0), distributed by arXiv under the Creative Commons Attribution 4.0 International licence, http://creativecommons.org/licenses/by/4.0/, as recorded in the arXiv licence metadata for this exact version and shown on the abstract page https://arxiv.org/abs/2512.17118v1. Full licence notice: \texttt{licenses/arxiv-2512.17118-CC-BY-4.0.txt}.\par\smallskip

\noindent\textit{Editorial scope.} Counted unit: the article's proposition 2.8 (source0.tex lines 229--231). The article's complete original proof of it is delivered with it (source0.tex lines 232--253, one \texttt{proof} environment of 22 lines). No mathematics is written, repaired or re-derived: the statement and the proof are the article's own lines, in the article's own notation, and no retained line is substituted or re-flowed.  No further source-local context is needed: every cross-reference of every retained line resolves inside the two delivered spans.  Nothing was omitted from any retained span, so the package carries no omission record.  No retained line contains a citation, so the wrapper carries no bibliography and no bibliography entry.  The retained lines contain the article's own diagram code, which the wrapper typesets with the standard tikz packages; no image file is delivered and the package declares no asset dependency.  Editorial formatting only, in addition: an AMS \texttt{amsart} preamble that restates the article's own declarations and macros used by the retained lines, namely the environments \texttt{proposition} on separate counters, together with the standard packages those lines need.  The retained environments print in this document as Proposition 1 in order of appearance, whereas the article prints the same items with the numbering of its own full document.  The complete native source of the article as distributed in the arXiv e-print for this exact version is bundled unchanged as \texttt{sources/191292/source0.tex}.
\begin{proposition}\label{2.8}
Let $F$ be a field and $x,y\in F^{*}$ such that $x^{a}=y^{b}$ only when $(a,b)=(0,0)\in\mathbb{Z}^{2}$, then $x\wedge y\neq0$ in $\wedge^{2}_{\mathbb{Z}}(F^{*})$.    
\end{proposition}

\begin{proof}
    Let $G:=<x,y>$, the subgroup of $F^{*}$ generated by $x$ and $y$. So, $G$ is a free abelian group of rank 2. Now, consider the following commutative diagram induced by the inclusion $i:G\rightarrow F^{*}$:

% https://q.uiver.app/#q=WzAsNixbMCwwLCJcXHdlZGdlXnsyfV97XFxtYXRoYmJ7Wn19RyJdLFsyLDAsIlxcd2VkZ2VeezJ9X3tcXG1hdGhiYntafX1GXnsqfSJdLFswLDEsIlxcd2VkZ2VeezJ9X3tcXG1hdGhiYntafX1HXFxiaWdvdGltZXNfe1xcbWF0aGJie1p9fVxcbWF0aGJie1F9Il0sWzIsMSwiXFx3ZWRnZV57Mn1fe1xcbWF0aGJie1p9fUZeeyp9XFxiaWdvdGltZXNfe1xcbWF0aGJie1p9fVxcbWF0aGJie1F9Il0sWzAsMiwiXFx3ZWRnZV57Mn1fe1xcbWF0aGJie1F9fShHXFxvdGltZXNfe1xcbWF0aGJie1p9fVxcbWF0aGJie1F9KSJdLFsyLDIsIlxcd2VkZ2VeezJ9X3tcXG1hdGhiYntRfX0oRl57Kn1cXG90aW1lc197XFxtYXRoYmJ7Wn19XFxtYXRoYmJ7UX0pIl0sWzAsMSwiXFx3ZWRnZV57Mn1fe1xcbWF0aGJie1p9fShpKSIsMl0sWzIsM10sWzAsMl0sWzEsM10sWzIsNCwiXFxjb25nIiwyXSxbMyw1LCJcXGNvbmciXSxbNCw1LCIiLDEseyJzdHlsZSI6eyJ0YWlsIjp7Im5hbWUiOiJob29rIiwic2lkZSI6InRvcCJ9fX1dXQ==
\[\begin{tikzcd}
	{\wedge^{2}_{\mathbb{Z}}G} && {\wedge^{2}_{\mathbb{Z}}F^{*}} \\
	{\wedge^{2}_{\mathbb{Z}}G\bigotimes_{\mathbb{Z}}\mathbb{Q}} && {\wedge^{2}_{\mathbb{Z}}F^{*}\bigotimes_{\mathbb{Z}}\mathbb{Q}} \\
	{\wedge^{2}_{\mathbb{Q}}(G\otimes_{\mathbb{Z}}\mathbb{Q})} && {\wedge^{2}_{\mathbb{Q}}(F^{*}\otimes_{\mathbb{Z}}\mathbb{Q})}
	\arrow["{\wedge^{2}_{\mathbb{Z}}(i)}"', from=1-1, to=1-3]
	\arrow[from=1-1, to=2-1]
	\arrow[from=1-3, to=2-3]
	\arrow[from=2-1, to=2-3]
	\arrow["\cong"', from=2-1, to=3-1]
	\arrow["\cong", from=2-3, to=3-3]
	\arrow[hook, from=3-1, to=3-3]
\end{tikzcd}\]

Now, the bottom arrow is injective as $G\otimes_{\mathbb{Z}}\mathbb{Q}$ is $\mathbb{Q}$-subspace of the vector space $F^{*}\otimes_{\mathbb{Z}}\mathbb{Q}$. Also, note that the one dimensional $\mathbb{Q}$-vector space $\wedge^{2}_{\mathbb{Q}}(G\otimes_{\mathbb{Z}}\mathbb{Q})$ is spanned by $(x\otimes_{\mathbb{Z}}1)\wedge_{\mathbb{Q}} (y\otimes_{\mathbb{Z}}1)$. If $x\wedge y=0$ in $\wedge^{2}_{\mathbb{Z}}F^{*}$, $(x\otimes_{\mathbb{Z}}1)\wedge_{\mathbb{Q}} (y\otimes_{\mathbb{Z}}1)=0$ in $\wedge^{2}_{\mathbb{Q}}(F^{*}\otimes_{\mathbb{Z}}\mathbb{Q})$, a contradiction to the bottom map being injective.


    
\end{proof}

\end{document}
